\section{Dense graphs with antitone colour degrees}\label{sec:dense}

The second global degree condition uses Tutte's factor theorem directly
on a doubled graph, where parallel edges are allowed by that theorem.
This is separate from the simple-host theorem used in
Section~\ref{sec:highdegree}.

\begin{theorem}[D3]\label{thm:dense}
Let $G$ be simple with $n$ vertices and
\[
 8\delta(G)\ge7n+6.
\]
Then $\doub G$ has a two-colouring in which all vertex red degrees
are pairwise distinct and all vertex blue degrees are pairwise distinct.
\end{theorem}
\begin{proof}
Put $q=n-1-\delta(G)$, so the condition is $n\ge8q+14$.
Order $v_0,\ldots,v_{n-1}$ by nonincreasing ordinary degree and put
$L=\lfloor(n-1)/2\rfloor$. Set $f(v_i)=L+i$, except that if the
sum is odd, increase only $f(v_{n-1})$ by one. The increment occurs
exactly when $n\equiv2\pmod4$; only evenness is needed below.
Their sum is even and the targets are strictly increasing. Their
physical range follows explicitly from
\[
 0\le f(v_i)\le L+n\le2(n-1-q)\le2d_G(v_i),
\]
where $L\le(n-1)/2$ and $n\ge8q+14$ imply the middle bound.
If $\doub G$ has an $f$-factor, use it for red.
Then $B(v_i)=2d_G(v_i)-f(v_i)$ is strictly decreasing, proving the
asserted strong conclusion.

We verify the factor inequalities. For disjoint $S,T$, put
$s=|S|,t=|T|,h=s+t,r=n-h$. Their Tutte left side is
\[
 D=f(S)+2\sum_{v\in T}d_{G-S}(v)-f(T).
\]
The odd-component term is the number $o_f(S,T)$ of components of
$G-S-T$ with odd $f$ sum, since all boundary edges in $\doub G$
have even multiplicity. Independently of their positions in the order,
\begin{align*}
 f(S)&\ge sL+s(s-1)/2,\\
 f(T)&\le t(L+n-1)-t(t-1)/2+1,\\
 d_{G-S}(v)&\ge n-1-q-s.
\end{align*}
The last lower bound remains valid when negative. Writing
$L=(n-1)/2-\theta$, $0\le\theta\le1/2$, gives
\begin{align*}
 D&\ge (s-t)L+\frac{s(s-1)+t(t-1)}2
       +t(n-1-2q-2s)-1\\
  &\ge h\left(\frac{n-6}{4}-2q\right)-1
   \ge2h-1.
\end{align*}
For the middle inequality use
$s^2+t^2-4st\ge-h^2/2\ge-nh/2$,
$-\theta(s-t)\ge-h/2$, and $t\le h$.

If $h=0$, the graph is connected: each component would contain at
least $n-q$ vertices and $n>2q$. Its total $f$ sum is even, so
$o_f=0=D$. If $h\ge1$ and $r>2q$, the remaining graph is connected
by the same argument, now with each component of size at least
$r-q>r/2$. Thus $o_f\le1\le2h-1$. If $r\le2q$, then $o_f\le r$
and $h=n-r\ge n-2q\ge6q+14\ge r$, so again $D\ge o_f$.
All factor inequalities hold. Tutte's theorem~\cite{Tutte54} supplies
the red factor.
\end{proof}

No optimality of the constant is claimed. The construction may use
all three weights $0,1,2$, unlike Theorem~\ref{thm:chromatic}.
Searching for a particular factor is an implementation task; a search
timeout is neither a failed existence proof nor a counterexample.

\section{Dense forbidden targets and sequential absorption}\label{sec:denselists}

Let $H$ have order $n$ and $q=n-1-\delta(H)$. At each vertex fix
physical external contributions $A_v,B_v$ and a finite set $F_v$ of
forbidden \emph{final red degrees}. Put
\[
 D_v=d_H(v)+(A_v+B_v)/2,\quad
 \alpha=\min_v A_v,\quad a=\max_v A_v-\alpha.
\]
Partition the vertices into equal-$D_v$ classes $C_1,\ldots,C_k$ in
strictly decreasing $D$ order and put
\[
 b_i=\max_{v\in C_i}|F_v|,\quad
 E=\sum_i b_i+\min_i b_i+2,\quad W=\max\{a,E\}.
\]

\begin{theorem}[D5: forbidden-target dense interface]\label{thm:denselists}
If $n\ge8q+4W+10$, one can colour $\doub H$ so that every final
red degree avoids $F_v$, all final red degrees are pairwise distinct,
and all final blue degrees are pairwise distinct. The external
contributions are preserved.
\end{theorem}
\begin{proof}
Let $L=\lfloor(n-1)/2\rfloor$. Starting at $\alpha+L$, allocate
pairwise disjoint consecutive integer blocks in the specified class
order. The block for $C_i$ has length $|C_i|+b_i$, except that a
class of minimum $b_i$ receives an additional $b_i+2$ places.
Together these blocks form one interval of length $n+E$.
In each ordinary class assign distinct targets avoiding the
individual forbidden sets by greedy choice; the length exceeds the
number of assigned predecessors plus any forbidden set.

In the enlarged minimum class, designate a vertex $z$. Its interval
has at least $2b_i+3$ places, containing at least $b_i+1$ integers
of each parity. Reserve one allowed even and one allowed odd value
for $z$. Assign the other vertices distinct allowed values while
avoiding both reserved places. Writing $n_i=|C_i|$, at the last
such step at least
$(n_i+2b_i+2)-2-b_i-(n_i-2)=b_i+2$ places remain after all
exclusions; if $n_i=1$ there is no such step. Choose the parity at $z$ that makes
$\sum_v(R_v-A_v)$ even. The resulting final red targets are distinct.
Order vertices inside each degree class by increasing chosen target.
In the resulting global order $v_0,\ldots,v_{n-1}$, $D_v$ is
nonincreasing, $R_v$ is strictly increasing, and therefore
$2D_v-R_v$ is strictly decreasing.

The $n$ selected places in an interval of length $n+E$ have ordered
indices $j_i$ with $i\le j_i\le i+E$. For the internal target
$f(v_i)=R_{v_i}-A_{v_i}$ we obtain
\begin{equation}\label{eq:densetargetband}
 L+i-a\le f(v_i)\le L+i+E.
\end{equation}
The assumption gives $L-a\ge0$ and
$L+n-1+E\le2(n-1-q)\le2d_H(v)$, so these are physical internal
targets. Their total is even by the reserved parity choice.

For disjoint $S,T$ with $s,t,h,r$ as before,
\begin{align*}
 f(S)&\ge sL+s(s-1)/2-as,\\
 f(T)&\le t(L+n-1)-t(t-1)/2+Et.
\end{align*}
Using $as+Et\le Wh$ and the same degree estimate as in the dense
proof yields, without its extra final parity increment,
\[
 f(S)+2\sum_{v\in T}d_{H-S}(v)-f(T)
 \ge h\left(\frac{n-6}{4}-2q-W\right)\ge h.
\]
For $h=0$ the connected graph has even total target, so its odd term
is zero. For $h\ge1,r>2q$ the remaining graph is connected and has
odd term at most one. For $r\le2q$ it has at most $r$ odd
components, while $h\ge n-2q\ge6q+4W+10\ge r$.
Thus every factor inequality holds on $\doub H$. The red factor
realises $f$, and adding the fixed external counts realises the
chosen final $R$ and $2D-R$. Their order and forbidden-value choices
prove the theorem.
\end{proof}

Theorem~\ref{thm:dense} remains useful: with no external data or
forbidden values it needs $n\ge8q+14$, whereas the general theorem
has $E=2$ and needs $n\ge8q+18$. It is therefore not removed as
a supposedly weaker duplicate.

\subsection{Preserving arbitrary fixed boundary colours}

Let $H=G[W]$ be induced, all outside and boundary weights fixed,
and all outside--outside edges legal using their actual counts.
For each fixed edge $vu$, $v\in W,u\notin W$, put in $F_v$
the value $R(u)$ if $x_{vu}>0$, and the value
$2d_G(v)-B(u)$ if $x_{vu}<2$. Let $A_v,B_v$ be its total fixed
incident contributions. A final red degree avoiding $F_v$ makes
every boundary comparison legal, while $D_v=d_G(v)$.
Theorem~\ref{thm:denselists} therefore gives a preserved-colour
extension under its stated bound.

If each region vertex has at most $t$ outside edges, then
$a\le2t$, $b_i\le2t$, and the possible total degrees lie between
$n-1-q$ and $n-1+t$. Thus $k\le q+t+1$ and
$W\le2t(q+t+2)+2$. A uniform sufficient condition is
\begin{equation}\label{eq:denseuniform}
 n\ge8q+8t(q+t+2)+18.
\end{equation}
The bound uses the per-vertex number of boundary edges, not their
total number or the degrees of their outside endpoints. For $q=0,t=1$
it is 42, and for $q=3,t=1$ it is 90. The actual class-dependent
bound can be smaller, for example in the supplied 34-vertex clique
witness with fixed mixed ports.

\begin{corollary}[Sequential dense absorption]\label{cor:sequential}
Let $W_1,\ldots,W_s$ be disjoint induced regions. Fix every edge
not internal to one region, including all inter-region edges. Assume
the low--low edges outside the regions are legal with their fixed
final counts. If each region satisfies~\eqref{eq:denseuniform}
with its internal complement degree $q_i$ and per-vertex outside
edge bound $t_i$, all region interiors can be filled in any order,
preserving all fixed edge weights. The sharper actual $W_i$
condition of Theorem~\ref{thm:denselists} may be used instead at
each step.
\end{corollary}
\begin{proof}
At a region include all fixed outside incidences in its offsets,
including those toward unprocessed regions. Their contributions are
known even if the other endpoints' final degrees are not yet known.
Forbid conflicts with low vertices and already processed regions,
whose degrees are final. Defer comparisons toward future regions.
The dense theorem fills the present interior, verifying all its
internal and earlier boundary edges. Later fillings cannot alter
these degrees or the fixed boundary weights; they enforce the
deferred comparisons at their own steps. Induction over the chosen
order verifies every edge. The uniform bound allows the worst case
of at most $2t_i$ forbidden values per vertex at every step.
\end{proof}

The placeholder-leaf preprocessing from
Section~\ref{sec:crtregions} supplies low-side colours when its
components are bipartite or cacti, with the singleton-port exception
handled explicitly. Unlike the CRT region theorem, sequential dense
absorption requires no proper colouring of the entire union of high
regions. Neither theorem asserts that an arbitrary graph admits a
partition meeting its degree hypotheses.
